\section{Einleitung}\label{einleitung}

Die Verbreitung von Internet of Things (\iot{})-Geraeten in privaten, industriellen und
kritischen Infrastrukturen vergroessert die Angriffsflaeche moderner Netzwerke
kontinuierlich~\cite{hamedani2025iotsurvey}.
Hamedani et al.\ \cite{hamedani2025iotsurvey} zeigen in einem umfassenden Survey ueber
ML-basierte Sicherheitsloesungen fuer \iot{}, dass die Forschung in diesem Feld zwischen 2020
und 2024 stark gewachsen ist und gleichzeitig methodische Luecken bestehen, insbesondere in
der Bewertung und der Interpretierbarkeit der eingesetzten Modelle.
Kikissagbe und Adda~\cite{kikissagbe2024review} kategorisieren in ihrem systematischen Review
supervised, unsupervised und hybride ML-Ansaetze fuer \iot{}-Intrusion Detection Systeme
(\ids{}) und ordnen baumbasierte Ensemble-Methoden als etablierten Forschungsstand ein.

Eine zentrale empirische Studie auf diesem Forschungsfeld ist die Arbeit von Raturi et al.\
\cite{raturi2026exploratory}, die mehrere ML-Algorithmen auf dem
CICIoT2023-Datensatz~\cite{neto2023ciciot2023} evaluiert.
\gb{} erzielt in der Mehrklassenklassifikation die hoechste \ba{} unter den getesteten
Modellen~\cite{raturi2026exploratory}.
Die Autoren benennen jedoch drei methodische Luecken: die fehlende Erklaerbarkeit der
Modellentscheidungen, die fehlende Analyse der Fehlerfortpflanzung im zweistufigen
Klassifikationsansatz und das Fehlen von Latenzmessungen~\cite{raturi2026exploratory}.
Die vorliegende Arbeit adressiert die erste dieser Luecken.

Die Wahl von SHapley Additive Explanations (\shap{}) als Erklaerbarkeitsmethode ist
wissenschaftlich begruendet.
Ogunseyi et al.\ \cite{ogunseyi2026xai_review} analysieren in einem PRISMA-basierten
systematischen Review 129 Studien aus den Jahren 2018 bis 2025 zur erklaerbaren Kuenstlichen
Intelligenz (XAI) in \ids{} fuer \iot{}-Netzwerke und zeigen, dass \shap{} in der
ueberwiegenden Mehrheit der untersuchten Studien eingesetzt wird, ohne die
Erkennungsgenauigkeit der zugrundeliegenden Modelle zu beeintraechtigen.
Mohale und Obagbuwa~\cite{mohale2025xai} demonstrieren am Beispiel mehrerer ML-Klassifikatoren
einschliesslich XGBoost, wie globale und lokale \shap{}-Analysen in \ids{}-Experimenten
umgesetzt werden koennen.
Hermosilla, Berrios und Allende-Cid~\cite{hermosilla2025xai_forensic} ergaenzen diesen
methodischen Rahmen um quantitative Bewertungsmetriken fuer die Qualitaet der Erklaerungen,
darunter Fidelitaet, Konsistenz und Jaccard-Aehnlichkeit.

Die vorliegende Arbeit verbindet diese methodischen Bausteine.
Sie entwickelt einen \gb{}-Klassifikator auf dem CICIoT2023-Datensatz, evaluiert ihn mit der
\ba{} als primaerer Metrik und integriert \shap{} als post-hoc-Erklaerungskomponente.
Eine Ablation Study untersucht zusaetzlich den Einfluss der Principal Component Analysis
(\pca{}) auf Modellleistung und Erklaerungen.

\textbf{Struktur der Arbeit.}
Kapitel~\ref{related_work} ordnet die zehn verwendeten Quellen thematisch in vier Cluster
ein und arbeitet die Forschungsluecke heraus.
Kapitel~\ref{methodik} beschreibt Datenbasis, Vorverarbeitung, Klassifikationsarchitektur,
Evaluationsmetriken und \shap{}-Integration.
Kapitel~\ref{evaluation} dokumentiert die Ergebnisse beider Preprocessing-Pipelines.
Kapitel~\ref{diskussion} diskutiert die Ergebnisse im Vergleich zur bestehenden Literatur.
Kapitel~\ref{fazit} schliesst mit Fazit und Ausblick.
